\Cref{tab:2-position} 汇总上述图索引、压缩搜索和自适应计算量方法的设计权衡。对 JHQ，需要研究的是两级表示如何影响 GPU 构建、候选选择与工作负载级精排。第 4--6 章展开 JHQ-GPU 的方法；第 7 章使用第 3 章规定的实测基线与对照进行评估。

每维九位的 IVF-RaBitQ 与 JHQ 每维一位主码加八位残差的名义编码载荷相同。性能前沿也包含每维四位和八位的 IVF-RaBitQ 配置；总显存用量单独评估。

\begin{table}[!htbp]
\centering
\setstretch{1.0}
\caption[相关方法与 JHQ-GPU 设计问题]{相关方法及其引出的 JHQ-GPU 设计问题。}
\label{tab:2-position}
\footnotesize\setlength{\tabcolsep}{4pt}
\begin{tabularx}{\linewidth}{@{}>{\raggedright\arraybackslash}p{.20\linewidth}*{2}{>{\raggedright\arraybackslash}X}@{}}
\toprule
方案 & 主要权衡 & JHQ-GPU 的问题 \\
\midrule
IVF-PQ / OPQ & 码本以降低重构误差为训练目标；子空间内部坐标可能耦合 & 主码的坐标结构如何简化 GPU 评分？ \\
Fast Scan / Quicker ADC & 小表在硬件和精度约束下加速查找 & 拆表能否抵消额外 GPU 查表？ \\
SONG / CAGRA & 图导航减少评分，但增加图存储和依赖访问 & 压缩扫描在哪些区间能与常驻图搜索竞争？ \\
BANG & 主机存储扩大容量，但增加通信 & 单 GPU 常驻完整压缩索引能实现什么？ \\
IVF-RaBitQ & 多位编码以存储和 probe 工作量换取距离估计质量 & 相同召回率下，JHQ 两阶段搜索与 IVF-RaBitQ 相比表现如何？ \\
自适应计算量 & 质量控制可能需要标签、预测或校准 & 工作负载一致性能否以可接受的质量和成本选择有效预算？ \\
\bottomrule
\end{tabularx}
\end{table}
